\section[Reflectionless limits and normal coordinates]{Reflectionless limits and compatible normal coordinates}
\label{sec:planar-limits}

We first compare two planes using points near both of them. This gives one
reference plane that controls the errors over a finite range of dyadic radii.
We then identify the
measure obtained when those errors vanish.  Three assertions enter that
identification separately: the limit is supported in a plane, its support
fills the plane, and its density on that plane is constant.

\subsection{Comparison from measured errors}

\begin{lemma}[Comparison of fitting planes]\label{lem:plane-comparison}
Fix $1\le n\le d$, $G\ge0$, and $c>0$.  There are constants
$A>0$ and $0<\varepsilon_0\le1$, depending only on $n,d,G,c$, with the
following property.  Suppose that $\sigma$ has global upper $n$-growth $G$,
$R>0$, and $\sigma(B(0,R))\ge cR^n$.  If $S,W$ are affine $n$-planes and
$0<\eta\le\varepsilon_0R$ satisfy
\begin{equation}\label{eq:joint-plane-error}
 \int_{B(0,R)}\bigl(\dist(x,S)^2+\dist(x,W)^2\bigr)\,d\sigma(x)
 \le \frac c2\eta^2R^n,
\end{equation}
then
\begin{align}
 \dist(x,W)&\le A\eta(1+|x|/R) &&(x\in S),\label{eq:plane-comparison}\\
 |\dist(x,S)-\dist(x,W)|&\le A\eta(3+|x|/R)
       &&(x\in\R^d).\label{eq:distance-comparison}
\end{align}
The first conclusion also holds with $S$ and $W$ interchanged.
\end{lemma}

\begin{proof}
Every integral in \eqref{eq:joint-plane-error} is finite, since its integrand
is bounded on the ball and upper growth implies local finiteness.  Markov's
inequality shows that
\[
 F=\{x\in B(0,R):\dist(x,S)^2+\dist(x,W)^2<\eta^2\}
 \quad\hbox{satisfies}\quad \sigma(F)\ge(c/2)R^n.
\]
In particular, every point of $F$ is within $\eta$ of both planes.

Here is the quantitative nonconcentration needed to choose points in $F$.
If $V$ is an affine $k$-plane, $k<n$, and $0<\tau\le1$, then
\begin{equation}\label{eq:lower-dimensional-tube}
 \sigma\{x\in\overline B(0,R):\dist(x,V)<\tau R\}
 \le G2^n3^k\tau R^n.
\end{equation}
Indeed, projection onto $V$, centered at the projection of $0$, sends
$\overline B(0,R)$ into its intrinsic closed $R$-ball.  A maximal
$\tau R$-separated subset of this intrinsic ball has at most
$(1+2/\tau)^k$ points: its disjoint open $\tau R/2$-balls lie in the
$(R+\tau R/2)$-ball, and comparison of their $k$-dimensional volumes gives
this bound.  Maximality gives a $\tau R$-net.  Ambient balls of radius
$2\tau R$ about its points cover the indicated tube.  Upper growth gives
$G2^n3^k\tau^{n-k}R^n$, which implies
\eqref{eq:lower-dimensional-tube}.

Choose
\[
 \tau=\min\{1,c/[4(G2^n3^n+1)]\}.
\]
The tube mass in \eqref{eq:lower-dimensional-tube} is strictly less than
$(c/2)R^n$.  Thus we can successively choose
$p_0,\ldots,p_n\in F\cap\supp\sigma$ such that, for $1\le j\le n$,
\[
 \dist\bigl(p_j,p_0+\operatorname{span}
              \{p_i-p_0:1\le i<j\}\bigr)\ge\tau R.
\]
At each step the affine span has dimension less than $n$, and $\supp\sigma$
has full measure.  Put $v_j=p_j-p_0$.  Then $|v_j|\le2R$.  These vectors
satisfy
\begin{equation}\label{eq:anchor-conditioning}
 \sum_{j=1}^n|t_j|\le \frac KR\left|\sum_{j=1}^nt_jv_j\right|,
\end{equation}
where $K$ depends only on $n,\tau$.  For completeness, define
$K_0=0$ and $K_{j+1}=K_j(1+2/\tau)+1/\tau$.
Projection perpendicular to the span of the first $j$ vectors bounds the
last coefficient by the norm of the complete linear combination divided
by $\tau R$.  Subtracting its last term enlarges that norm by at most
$1+2/\tau$.  Induction proves \eqref{eq:anchor-conditioning} with $K=K_n$.
It also proves independence of the $v_j$.

Let $a=\pi_Sp_0$ and $w_j=\pi_Sp_j-a$, where $\pi_S$ is orthogonal
projection.  We have $|w_j-v_j|\le2\eta$.  If
$\eta/R\le1/[4(K+1)]$, \eqref{eq:anchor-conditioning} and absorption give
\[
 \sum|t_j|\le\frac{2K}{R}\left|\sum t_jw_j\right|.
\]
The $n$ vectors $w_j$ consequently form a basis of the direction of $S$.
Every projected anchor $a,a+w_1,\ldots,a+w_n$ lies within $2\eta$ of $W$.
The orthogonal normal-coordinate map for $W$ is affine and has norm
$\dist(\cdot,W)$.  Applying this affine map to
$x=a+\sum t_jw_j$ gives
\[
 \dist(x,W)\le2\eta+4\eta\sum|t_j|
 \le2\eta+\frac{8K\eta}{R}|x-a|.
\]
Since $|a|\le R+\eta\le2R$, this proves
\eqref{eq:plane-comparison}, for example with
$A=2(1+8K)$ and
$\varepsilon_0=\min\{1,1/[4(K+1)]\}$.
All choices preceded $\sigma,R,S,W,\eta$.

Finally, for arbitrary $x$, let $p=\pi_Sx$.  Since
$|\pi_S0|\le|a|\le2R$ and the linear part of $\pi_S$ is a contraction,
$|p|\le|x|+2R$.  Hence
\[
 \dist(x,W)\le|x-p|+\dist(p,W)
 \le\dist(x,S)+A\eta(3+|x|/R).
\]
Interchanging the planes proves \eqref{eq:distance-comparison}.
\end{proof}

\begin{lemma}[One reference plane through a finite horizon]
\label{lem:fixed-plane-moments}
Fix $0<\gamma<1/2$ and the constants $n,d,G,c$ of
Lemma~\ref{lem:plane-comparison}.  There are $\kappa>0$ and $M<\infty$,
depending only on these data and $\gamma$, as follows.  Suppose $K\ge0$ is
an integer, $\delta>0$, $\delta2^{\gamma K}\le\kappa$, $\sigma$ has
global upper growth $G$, and, for $0\le k\le K$,
\[
 \sigma(B(0,2^k))\ge c2^{kn},
 \qquad e_\sigma(0,2^k)\le\delta2^{\gamma k}.
\]
Then an affine $n$-plane $S$ satisfies
\begin{equation}\label{eq:fixed-plane-moments}
 \int_{B(0,2^k)}\dist(x,S)^2\,d\sigma(x)
 \le M\delta^2 2^{k(n+2+2\gamma)}
 \qquad(0\le k\le K).
\end{equation}
\end{lemma}

\begin{proof}
Write $b=2^\gamma>1$ and $R_i=2^i$.  The definition of the infimum in
$e_\sigma$ and the positivity of $\delta b^i$ give affine planes $P_i$,
$0\le i\le K$, with
\[
 \int_{B(0,R_i)}\dist(x,P_i)^2\,d\sigma
 \le2(\delta b^i)^2R_i^{n+2}.
\]
No attainment of a least-squares infimum is required.  For $0\le i<K$,
restricting the error for $P_{i+1}$ to $B(0,R_i)$ shows that the sum of neighboring errors
is at most
$D(\delta b^i)^2R_i^{n+2}$, where
$D=2+2\cdot2^{n+2}b^2$.
Choose $L>0$ with $D\le L^2c/2$, and put
$\eta_i=L\delta R_i b^i$.  Taking
$\kappa=\varepsilon_0/L$ ensures $\eta_i\le\varepsilon_0R_i$ through
the horizon.  Lemma~\ref{lem:plane-comparison} gives, at the same arbitrary
point $x$ in each step,
\[
 \dist(x,P_i)\le\dist(x,P_{i+1})
       +AL\delta b^i(3R_i+|x|).
\]
For $|x|\le R_k$, the sum of these errors for $i<k$ is bounded by
$B\delta R_kb^k$, where
\[
 B=AL\left(\frac3{2b-1}+\frac1{b-1}\right).
\]
Indeed, sum the geometric series with ratios $2b$ and $b$ separately.
Consequently
$\dist(x,P_0)\le\dist(x,P_k)+B\delta R_kb^k$ on that ball.
Squaring, integrating, and using upper growth proves
\eqref{eq:fixed-plane-moments} with $S=P_0$ and $M=4+2GB^2$.
\end{proof}

\subsection{Rigidity on the limiting plane}

We shall use the symmetric fractional operator
\begin{equation}\label{eq:fractional-operator-definition}
 \mathcal Lg(x)=\int_{\R^n}
 \frac{2g(x)-g(x+z)-g(x-z)}{2|z|^{n+1}}\,dz,
 \qquad g\in C_c^\infty(\R^n).
\end{equation}
This integral is absolutely convergent: the numerator is $O(|z|^2)$ near
zero and is bounded at infinity. Lemma~\ref{lem:fractional-rigidity}
proves that a measurable function $v:\R^n\to\R$ satisfying
$\int|v(x)|(1+|x|)^{-n-2}\,dx<\infty$ and
$\int v\mathcal Lg=0$ for every compact smooth mean-zero test is affine
almost everywhere. If $v$ is essentially bounded, it is constant
almost everywhere.
That lemma concerns functions on $\R^n$ and is proved independently
of the limiting measures constructed here.

\begin{lemma}[Rigidity of a planar reflectionless measure]
\label{lem:planar-rigidity}
Let $1\le n\le d$, let $L$ be an affine $n$-plane, and let $\sigma$ be a
nonzero Radon measure on $\mathbb R^d$ supported in $L$.  Suppose that
$c,G>0$ and
\[
 \sigma(B(x,r))\le Gr^n\quad(x\in\mathbb R^d,\ r>0),
 \qquad
 \sigma(B(x,r))\ge cr^n\quad(x\in\supp\sigma,\ r>0).
\]
Suppose also that $\mathcal R_\sigma(\varphi)=0$ for every compactly
supported Lipschitz function $\varphi$ with $\int\varphi\,d\sigma=0$.
Assume the hard Riesz truncations are uniformly bounded on $L^2(\sigma)$.
Then
\[
 \supp\sigma=L,
 \qquad \sigma=\theta\,\mathcal H^n\!\restriction L
 \quad\text{for some }\theta>0.
\]
The same conclusion holds if the operator assumption on $\sigma$ is
replaced by the following source assumption: $\sigma_j$ are Radon
measures on $\mathbb R^d$ converging to $\sigma$ against compact continuous
tests, with a common global upper $n$-growth bound and a common uniform
hard-truncation $L^2$ bound.  With the Euclidean normalization of
$\mathcal H^n$, the resulting constant satisfies
$c/\omega_n\le\theta\le G/\omega_n$, where
$\omega_n=|B_{\mathbb R^n}(0,1)|$.
\end{lemma}

\begin{proof}
An affine isometry identifies $L$ with $\mathbb R^n$.  Under this
identification write $\mu$ for $\sigma$.  The kernel and its scalar
projections commute with this isometry.  Reflectionlessness passes to
$\mu$ with its actual measure: extend an intrinsic compact Lipschitz test
by composing with orthogonal projection and multiplying by an ambient
compact cutoff equal to one on its trace support.  Its integral is
unchanged.  Choose a localization ball containing the extension and use
the localization independence of the mean-zero pairing from
Lemma~\ref{lem:pairings}.  Thus no density or full-support assertion is
needed to make this identification.

\emph{Full support.}
We first derive an annular bound using only growth and
reflectionlessness.  For a test supported in $\overline B(a,R/2)$, denote by
$P_{a,R}(\varphi)$ the vector expression
\begin{align*}
 &\frac12\iint_{B(a,R)^2}
 K_n(x-y)\bigl(\varphi(x)-\varphi(y)\bigr)\,d\mu(x)d\mu(y)\\
 &\qquad+\int_{B(a,R)}\varphi(x)
 \int_{B(a,R)^c}\bigl[K_n(x-y)-K_n(a-y)\bigr]\,d\mu(y)d\mu(x).
\end{align*}
Both terms are absolutely integrable by the Lipschitz cancellation near
the diagonal and the extra power of decay in the exterior difference.
This expression is defined without a mean-zero assumption; on mean-zero
tests it equals $\mathcal R_\mu$.  Splitting the larger ball into the
smaller ball and its complement, and using oddness, gives
\begin{equation}\label{eq:planar-radius-difference}
 P_{a,R}(\varphi)-P_{a,r}(\varphi)
 =\left(\int\varphi\,d\mu\right)
   \int_{B(a,R)\setminus B(a,r)}K_n(a-y)\,d\mu(y)
\end{equation}
when $0<r\le R$ and $\supp\varphi\subset\overline B(a,r/2)$.
One may first perform the splitting with a finite outer cutoff; the
integrable exterior differences justify its removal.

Choose a nonnegative Lipschitz bump equal to one on $B(a,s)$ and zero
outside $B(a,2s)$, and divide it by its $\mu$-integral.  For $a\in\supp\mu$
the resulting function $\eta_s$ has integral one and satisfies
\[
 \|\eta_s\|_\infty\le (cs^n)^{-1},\qquad
 \Lip(\eta_s)\le (cs^{n+1})^{-1},\qquad
 |P_{a,4s}(\eta_s)|\le C_n(G^2/c+G)=:A.
\]
For the last estimate the interior term is at most a constant times
$(cs^{n+1})^{-1}(Gs^n)(Gs)$; the exterior term is at most a constant
times $Gs/s$, since $\int\eta_s\,d\mu=1$.
Apply reflectionlessness to $\eta_{r/4}-\eta_{R/4}$ and then use
\eqref{eq:planar-radius-difference}.  It follows that
\[
 \left|\int_{B(a,R)\setminus B(a,r)}K_n(a-y)\,d\mu(y)\right|
 \le 2A\qquad(0<r\le R).
\]
Letting inner radii decrease strictly to $\varepsilon$ gives the same
bound on $\{y:\varepsilon<|y-a|<R\}$.  At fixed
$\varepsilon>0$ the kernel is integrable there, so this passage does not
require zero mass on spheres.

Suppose now that $z\notin\supp\mu$, and let $a$ be a nearest support
point to $z$.  Translate $a$ to zero and put $e=a-z\ne0$.
The nearest-point property gives
\[
 2\langle e,x\rangle+|x|^2\ge0\qquad(x\in\supp\mu).
\]
Define
$p_e(x)=\max(0,\langle e,x\rangle)/|x|^{n+1}$ for $x\ne0$,
and set $p_e(0)=0$.  Then
\[
 p_{-e}(x)\le\tfrac12|x|^{1-n},\qquad
 p_e(x)=\langle e,K_n(x)\rangle+p_{-e}(x).
\]
Upper growth makes the first majorant integrable on bounded balls: its
integral on $B(0,R)$ is at most $C_nGR$, as follows by dyadic annuli.
The signed annular bound therefore bounds the nonnegative truncated
integrals of $p_e$ on $B(0,1)$ uniformly.  Monotone convergence proves
$p_e\in L^1(\mu|_{B(0,1)})$.

Take a weakly convergent subsequence of the normalized blowups
$\mu_r(A)=r^{-n}\mu(rA)$ as $r\downarrow0$.
Upper growth supplies compactness, and the lower bound at zero makes the
limit $\tau$ nonzero.  These are the conclusions of
Lemma~\ref{lem:weak-compactness}, applied to the blowups.  The rescaled
quadratic inequality is
$2\langle e,x\rangle+r|x|^2\ge0$; hence
$\supp\tau\subset\{\langle e,x\rangle\ge0\}$.
For every fixed $T>0$,
\[
 \int_{B(0,T)}\max(0,\langle e,x\rangle)\,d\mu_r(x)
 \le T^{n+1}\int_{B(0,rT)}p_e(y)\,d\mu(y)\longrightarrow0.
\]
Testing with compact nonnegative cutoffs shows that $\tau$ is also
supported in $\{\langle e,x\rangle\le0\}$.  It is consequently supported
on a proper hyperplane in $\mathbb R^n$.  This is impossible for a
nonzero measure with upper $n$-growth: covering open sets by disjoint
dyadic cubes and enclosing each cube in a ball shows that such a measure
is bounded above by $C_nG'\,dx$, where $G'$ is its growth constant.
A proper hyperplane has Lebesgue measure zero.  This contradiction proves
$\supp\mu=\mathbb R^n$.

\emph{A positive bounded density.}
The same cube comparison gives $\mu=f\,dx$ with $f\le C_nG$ almost
everywhere.  Full support makes the lower ball bound valid at every
point.  Lebesgue differentiation then gives $f\ge c/\omega_n$ almost
everywhere.  Choose a measurable representative with
\begin{equation}\label{eq:planar-two-sided-density}
 0<c_1\le f(x)\le C_1<\infty\qquad(x\in\mathbb R^n).
\end{equation}
Changing $f$ on a null set has not changed $\mu$.

\emph{A local representative for the pairing.}
Fix a unit direction $u\in\mathbb R^n$, a center $a$, and radii
$0<H$ and $2H\le R$.  Put $B=B(a,R)$ and
\[
 T_{\varepsilon,B}1(x)=\int_B
 \frac{\langle u,x-y\rangle}{\max(\varepsilon,|x-y|)^{n+1}}\,d\mu(y).
\]
There is a constant $D_1$, independent of $B$ and $\varepsilon>0$, such
that
\begin{equation}\label{eq:planar-local-cap-bound}
 \|T_{\varepsilon,B}1\|_{L^2(\mu|_B)}
 \le D_1\mu(B)^{1/2}.
\end{equation}
For the direct operator hypothesis, subtract the hard truncation.
The error kernel is bounded by
$\varepsilon^{-n}\mathbf1_{\{|x-y|\le\varepsilon\}}$.
Its row and column integrals are at most $2^nG$, so Schur's inequality
gives \eqref{eq:planar-local-cap-bound} from the original operator bound.
Restricting the hard operator to a ball preserves that bound, by extending
the input by zero.

Here is also the justification under the source assumption.  Work first
in $\mathbb R^d$ on the ambient ball corresponding to $B$.  Its boundary
has $\sigma$-measure zero by \eqref{eq:planar-two-sided-density}.
Restricted source measures therefore converge weakly as finite measures.
The same Schur estimate gives uniformly bounded capped bilinear pairings
on every source restriction.  For fixed $\varepsilon>0$ these pairings
pass to the limit against two bounded continuous tests: the capped kernel
is bounded continuous, as are the squared tests giving their $L^2$
norms.  Take the input test to be one and the other test to be the capped
output for the limiting finite measure.  That output is bounded
continuous, by dominated convergence at this fixed cap.  The resulting
bilinear estimate gives its genuine $L^2$ bound
\eqref{eq:planar-local-cap-bound}.  Isometric pullback gives precisely the
displayed intrinsic transform.  This argument passes no singular kernel
through weak convergence.

Choose a subsequence $\varepsilon_j\downarrow0$ along which
$T_{\varepsilon_j,B}1$ converges weakly in $L^2(\mu|_B)$ to $v$.
This one subsequence works against every test in that Hilbert space.
Oddness gives, for bounded Lipschitz $\varphi$,
\[
 \int_B T_{\varepsilon,B}1\,\varphi\,d\mu
 =\frac12\iint_{B^2}
 \frac{\langle u,x-y\rangle(\varphi(x)-\varphi(y))}
      {\max(\varepsilon,|x-y|)^{n+1}}\,d\mu(x)d\mu(y).
\]
Removing the cap on the right changes the integral by at most
$C_nG\Lip(\varphi)\mu(B)\varepsilon$.  Thus $v$ represents the actual
interior singular pairing.  On $B(a,H)$ the function
\[
 F(x)=\int_{B(a,R)^c}
 \langle u,K_n(x-y)-K_n(a-y)\rangle\,d\mu(y)
\]
is well defined and bounded: the kernel difference is at most
$C_n|x-a||y-a|^{-n-1}$, whose exterior integral is at most
$C_nG|x-a|/R$.  Fubini now identifies the scalar component of
$\mathcal R_\mu(\varphi)$ with $\int(v+F)\varphi\,d\mu$ for mean-zero
compact Lipschitz tests supported in $B(a,H)$.

Choose a compact Lipschitz test $\eta$ supported in $B(a,H)$ with
$\int\eta\,d\mu=1$, using the positive mass of that ball, and put
$b=\int(v+F)\eta\,d\mu$.  For any smooth compact test $q$ in this smaller
ball, reflectionlessness applied to
$q-(\int q\,d\mu)\eta$ gives
$\int(v+F-b)q\,d\mu=0$.
Since $\mu=f\,dx$, the fundamental lemma for locally integrable functions
applied to $(v+F-b)f$ proves
\begin{equation}\label{eq:planar-local-constant}
 v+F=b\quad\text{$\mu$-almost everywhere on }B(a,H).
\end{equation}
This is a statement about the transform; constancy of $f$ has not yet
been used or proved.

\emph{An unweighted equation and density constancy.}
Let $g\in C_c^\infty(B(a,H))$ be real with $\int g\,dx=0$.
By \eqref{eq:planar-two-sided-density}, $g/f$ belongs to
$L^2(\mu|_B)$.  It is therefore an allowed test for the weak convergence,
even though it need not be Lipschitz.  Moreover, $\mu$ and Lebesgue measure
have the same null sets.  Testing with $g/f$ and using
\eqref{eq:planar-local-constant} yields
\begin{equation}\label{eq:planar-unweighted-local}
 \lim_j\left[\int g(x)T_{\varepsilon_j,B}1(x)\,dx
                   +\int g(x)F(x)\,dx\right]=0.
\end{equation}
The subsequence and $b$ were selected before $g$.

For a compact smooth function $g$, define the scalar test transform
\[
 R_ug(x)=-\frac12\int_{\mathbb R^n}
 \frac{\langle u,h\rangle[g(x+h)-g(x-h)]}{|h|^{n+1}}\,dh.
\]
The integral is absolutely convergent: the first difference is
$O(|h|)$ near zero, and the translated compact supports control infinity.
The analogous capped expression equals the capped transform of $g$ by
translation and reflection.  It converges pointwise to $R_ug$ and is
uniformly bounded on bounded $x$-sets.  If $\supp g\subset B(a,H)$ and
$\int g=0$, then for $|x-a|\ge2H$ and $\varepsilon\le H$ the cap is
inactive and mean-zero cancellation gives
\[
 |R_{u,\varepsilon}^{\rm cap}g(x)|
 \le C_n H\|g\|_1|x-a|^{-n-1}.
\]
Together these bounds form an integrable majorant independent of small
$\varepsilon$.  Since $f$ is bounded, dominated convergence applies to
$\int f R_{u,\varepsilon}^{\rm cap}g$.

For $\varepsilon\le H$, oddness and Fubini identify the expression in
brackets in \eqref{eq:planar-unweighted-local} with
$-\int f R_{u,\varepsilon}^{\rm cap}g\,dx$.
The local capped product is absolutely integrable.  For the exterior
term, use the centered kernel difference, which is absolutely integrable
against $|g(x)|\,dx\,d\mu(y)$; cancel the center term using $\int g=0$
inside the $x$-integral before integrating over the infinite exterior.
The open ball and its closed exterior are exact complements.  Thus no
separate divergent exterior integral or omitted boundary term occurs.
It follows that
\begin{equation}\label{eq:planar-density-riesz}
 \int_{\mathbb R^n}f(x)R_ug(x)\,dx=0
 \qquad\left(g\in C_c^\infty(\mathbb R^n),\ \int g=0\right).
\end{equation}
The ball may be chosen after the test, so this statement has no fixed
support restriction.

Let $u_1,\ldots,u_n$ be an orthonormal basis and
$\psi\in C_c^\infty(\mathbb R^n)$ be real.  Each
$g_i=\partial_{u_i}\psi$ has integral zero.  We claim
\begin{equation}\label{eq:planar-riesz-gradient}
 \sum_{i=1}^nR_{u_i}(\partial_{u_i}\psi)(x)=\mathcal L\psi(x).
\end{equation}
Indeed, set
\[
 Q_x(t)=\int\frac{\psi(x+th)+\psi(x-th)-2\psi(x)}{|h|^{n+1}}\,dh.
\]
Scaling gives $Q_x(t)=tQ_x(1)$ for $t>0$.  Differentiation under the
integral at $t=1$ is justified by second-derivative cancellation near
zero; for $1/2\le t\le2$ the derivative integrand vanishes for all
sufficiently large $|h|$ at fixed $x$.  Consequently
\[
 Q_x(1)=\int\frac{D\psi(x+h)h-D\psi(x-h)h}{|h|^{n+1}}\,dh.
\]
Expand $h$ in the orthonormal basis and use the definition of $R_u$ to
obtain \eqref{eq:planar-riesz-gradient}.  Every pairing with bounded $f$
is integrable by the preceding estimates.  Summing
\eqref{eq:planar-density-riesz} therefore gives
\[
 \int f\mathcal L\psi\,dx=0\qquad(\psi\in C_c^\infty(\mathbb R^n)).
\]
Real and imaginary parts give the same identity for complex tests.
Boundedness supplies $\int|f(x)|(1+|x|)^{-n-2}\,dx<\infty$.
Lemma~\ref{lem:fractional-rigidity} now makes $f$ affine almost everywhere.
A bounded affine function is constant: any nonzero linear part is
unbounded along a line, and continuity promotes an almost-everywhere
bound to an everywhere bound.  Thus $f=\theta$ almost everywhere, with
$\theta>0$ by \eqref{eq:planar-two-sided-density}.  The original ball
bounds give $c\le\theta\omega_n\le G$, completing the proof.
\end{proof}

\subsection{A common flat limit and its normal coordinates}

The passage from vanishing mean-zero pairings to a reflectionless weak
limit is part of the compactness method of Jaye and Nazarov
\cite{JayeNazarov2018I}. Here the common growth bounds control the
singular diagonal and exterior tails, as proved in
Lemma~\ref{lem:weak-compactness}. We combine that local argument with
the measured plane errors; we do not invoke a classification of
arbitrary higher-codimension reflectionless measures.

\begin{proposition}[The flat limit]\label{prop:flat-limit}
Fix $1\le n<d$, $c>0$, $0<G<\infty$, $0\le D_0<\infty$, and
$0<\gamma<1/2$.  Let $\mu_j$ be Radon measures on $\R^d$ with
$0\in E_j=\supp\mu_j$, global upper $n$-growth $G$, and
\[
 \mu_j(B(x,r))\ge cr^n
 \quad(x\in E_j,\ 0<r\le\diam E_j).
\]
Assume every hard-truncated Riesz transform on $L^2(\mu_j)$ has norm at
most $D_0$.  Let $\delta_j>0$, integers $K_j\ge0$, and $A_j,v_j>0$
satisfy
\begin{equation}\label{eq:flat-limit-horizons}
 \begin{gathered}
 \delta_j\longrightarrow0,\quad K_j\longrightarrow\infty,\quad
 \delta_j2^{\gamma K_j}\longrightarrow0,\\
 2^{K_j}\le\diam E_j,\quad A_j\longrightarrow\infty,\quad v_j\longrightarrow0.
 \end{gathered}
\end{equation}
Suppose
\begin{equation}\label{eq:flat-limit-source-excess}
 e_{\mu_j}(0,2^k)\le\delta_j2^{\gamma k}
 \qquad(0\le k\le K_j)
\end{equation}
and
\begin{equation}\label{eq:flat-limit-source-oscillation}
 |\mathcal R_{\mu_j}(\varphi)|\le v_j
 \quad\text{if }\ \Lip\varphi\le1,\quad
 \supp\varphi\Subset B(0,A_j),\quad\int\varphi\,d\mu_j=0.
\end{equation}
There is a subsequence, relabeled by $j$, with the following properties.
All the input sequences are restricted along this same subsequence when
we relabel them.
For each index there are an affine plane
$S_j=a_j+L_j(\R^n)$, tangent and normal linear isometries
\[
 L_j:\R^n\longrightarrow\R^d,
 \qquad N_j:\R^{d-n}\longrightarrow\R^d,
 \qquad L_jL_j^*+N_jN_j^*=I,
\]
such that $a_j$ is the point of $S_j$ nearest $0$,
$a_j\to0$, and $L_j\to L_\infty$ in operator norm, where $L_\infty$ is
an isometry.  Along this same subsequence,
\begin{equation}\label{eq:flat-limit-measure}
 \mu_j\rightharpoonup q(L_\infty)_\#\mathcal L^n,
 \qquad \frac{c}{4^n\omega_n}\le q\le\frac{2^nG}{\omega_n},
 \qquad \omega_n=\mathcal L^n(B(0,1)).
\end{equation}
Here convergence is against continuous compactly supported tests, and
$\mathcal L^n$ in \eqref{eq:flat-limit-measure} denotes Lebesgue measure.
The planes obey \eqref{eq:fixed-plane-moments} with a single constant
$M=M(n,d,c,G,\gamma)$.  In particular, writing $e_i$ for the standard
basis of $\R^{d-n}$ and
\begin{equation}\label{eq:flat-limit-normal-coordinate}
 u_{j,i}(x)=\langle N_je_i,x-a_j\rangle,
\end{equation}
we have
\begin{equation}\label{eq:flat-limit-normal-moments}
 \int_{B(0,2^k)}|u_{j,i}|^2\,d\mu_j
 \le M\delta_j^2 2^{k(n+2+2\gamma)}
 \qquad(0\le k\le K_j, 1\le i\le d-n).
\end{equation}
Moreover, on every bounded ball, the supports converge bilaterally to
$L_\infty(\R^n)$: each source support point approaches that plane uniformly,
and each point of the plane approaches $E_j$ uniformly.
\end{proposition}

\begin{proof}
After discarding finitely many indices, the smallness threshold in
Lemma~\ref{lem:fixed-plane-moments} holds.  That lemma gives $S_j$ and the
uniform moment bounds.  Each fixed-radius squared-distance integral tends
to zero: choose a larger fixed dyadic radius first, and then use
$K_j\to\infty$ and $\delta_j\to0$.

This integral conclusion also controls all support points in bounded
balls.  If $x_j\in E_j\cap B(0,R)$ and $\dist(x_j,S_j)\ge t>0$, the
$1$-Lipschitz distance function is at least $t/2$ on
$B(x_j,\min\{1,t/2\})$.  This ball lies in $B(0,R+1)$, and its radius
is admissible for all large $j$.  The lower growth bound therefore forces
a fixed positive lower bound on the squared-distance integral there,
a contradiction.  Thus
\begin{equation}\label{eq:moving-support-tube}
 \sup_{x\in E_j\cap B(0,R)}\dist(x,S_j)\longrightarrow0
 \quad\text{for every fixed }R.
\end{equation}
Because $0\in E_j$, the nearest points $a_j$ tend to $0$.
Choose orthonormal tangent and complementary normal frames $L_j,N_j$.
The compactness of the finite-dimensional set of tangent isometries gives
a subsequence with $L_j\to L_\infty$; norm preservation shows that the
limit is still an isometry.  The normal frames are retained along this
subsequence; their convergence is not required.

Lemma~\ref{lem:weak-compactness}, followed if necessary by a further
subsequence, gives a nonzero limit $\nu$ with $0\in\supp\nu$, hard
truncation bound $D_0$, global upper growth $2^nG$, and
\begin{equation}\label{eq:flat-limit-lower-growth}
 \nu(B(x,r))\ge4^{-n}cr^n
 \qquad(x\in\supp\nu, r>0).
\end{equation}
All frames and moment bounds are kept along the composition of the two
subsequences.  If $x\in\supp\nu$, positivity of compact tests in every
neighborhood of $x$ supplies nearby points of $E_j$ for all large $j$.
Together with \eqref{eq:moving-support-tube}, this gives
$\dist(x,S_j)\to0$.  On the other hand the affine projections
\[
 a_j+L_jL_j^*(x-a_j)\longrightarrow L_\infty L_\infty^*x.
\]
It follows that $x=L_\infty L_\infty^*x$.  At this stage we have proved
only $\supp\nu\subset L_\infty(\R^n)$.

We next verify that this same $\nu$ is reflectionless. Let $\varphi$ be a
compactly supported Lipschitz function with $\int\varphi\,d\nu=0$.
The pairing-convergence assertion of Lemma~\ref{lem:weak-compactness}
provides compact Lipschitz tests $\varphi_j$ of zero $\mu_j$-mean,
with uniformly bounded supports and Lipschitz constants, such that
\[
 \mathcal R_{\mu_j}(\varphi_j)\longrightarrow
 \mathcal R_\nu(\varphi).
\]
Choose $C\ge1$ bounding their Lipschitz constants. For all sufficiently
large $j$, their supports lie compactly inside $B(0,A_j)$, so
\eqref{eq:flat-limit-source-oscillation}, applied to $\varphi_j/C$,
gives $|\mathcal R_{\mu_j}(\varphi_j)|\le Cv_j\to0$.
Thus $\mathcal R_\nu(\varphi)=0$. This argument applies to every such
test on the measure subsequence already selected above.

Apply Lemma~\ref{lem:planar-rigidity} to this plane-supported limit, using
the inherited hard-truncation bound $D_0$.  It gives
$\nu=q(L_\infty)_\#\mathcal L^n$ with $q>0$.
Its upper growth and \eqref{eq:flat-limit-lower-growth}, evaluated on a
ball centered on the plane, give the two bounds for $q$ in
\eqref{eq:flat-limit-measure}.
Since $|u_{j,i}(x)|\le|N_j^*(x-a_j)|=\dist(x,S_j)$, the normal moment
bounds follow from those of $S_j$.

Finally, convergence of the affine projections is uniform on bounded
sets, so \eqref{eq:moving-support-tube} gives uniform closeness of source
support points to the limiting plane.  Conversely, if points $y_j$ of
that plane in a fixed bounded ball stayed a positive distance from $E_j$,
a subsequence would converge to a point $y$ of the plane.  A sufficiently
small ball about $y$ would then miss $E_j$ for all large $j$.  A nonnegative
compact test in that ball with positive integral against
$q(L_\infty)_\#\mathcal L^n$ contradicts weak convergence.  This proves
the stated bilateral convergence.
\end{proof}

In the later contradiction argument we take
$\gamma=1/4$, $\delta_j=2^{-j}$, $K_j=2j+3$, $A_j=8\cdot2^{K_j}$,
and $v_j=\delta_j^3$.  These choices satisfy
$\delta_j2^{\gamma K_j}\to0$ as well as
$\delta_j^{-1}2^{-K_j}\to0$.  The latter condition will control the
normalized exterior remainder; it was not needed to identify the
unnormalized planar measure.
